% Generated draft: requires mathematical and source review.
\documentclass[11pt]{article}
\usepackage[margin=25mm]{geometry}
\usepackage{fontspec}
\setmainfont{Latin Modern Roman}
\usepackage{amsmath,amssymb,mathtools,unicode-math}
\setmathfont{Latin Modern Math}
\usepackage{enumitem,needspace,xurl,hyperref}
\setlength{\parindent}{0pt}
\setlength{\parskip}{5pt}
\setlength{\emergencystretch}{3em}
\newcommand{\N}{\mathbb{N}}
\newcommand{\Z}{\mathbb{Z}}
\newcommand{\Q}{\mathbb{Q}}
\newcommand{\R}{\mathbb{R}}
\newcommand{\C}{\mathbb{C}}
\newcommand{\F}{\mathbb{F}}
\newcommand{\A}{\mathbb{A}}
\newcommand{\PP}{\mathbb{P}}
\newcommand{\E}{\mathbb{E}}
\newcommand{\eps}{\varepsilon}
\newcommand{\dd}{\,\mathrm{d}}
\begin{document}
\section*{UnsolvedMath ID: 11100022}
\textbf{Source problem code:} catalogue entry
\subsection{Definitions and mathematical statement}
For an ordinary vector bundle trivialized over \(\{U_i\}\), transition maps satisfy \(g_{ij}g_{jk}=g_{ik}\) on triple overlaps. Hovey proposes replacing this strict cocycle condition by coherent natural isomorphisms in the Kapranov--Voevodsky \(2\)-category of \(2\)-vector spaces, and asks whether this gives a useful construction of elliptic cohomology.

In a standard field-\(k\) model, a finite \(2\)-vector space is a \(k\)-linear category equivalent to a finite product of finite-dimensional vector-space categories; \(1\)-morphisms are linear functors and \(2\)-morphisms natural transformations. A precise descent proposal requires local objects, overlap equivalences \(g_{ij}\), and invertible \(2\)-morphisms
\[
 \alpha_{ijk}:g_{ij}\circ g_{jk}\Longrightarrow g_{ik}.
\]
On quadruple overlaps, the two composites from \(g_{ij}g_{jk}g_{kl}\) to \(g_{il}\) must agree, with associativity and unit coherence included. Merely allowing unrelated triple-overlap isomorphisms is not enough.

The mathematical program is to specify this descent theory and determine what classifying/cohomological construction it yields, and whether it realizes an intended elliptic cohomology theory. The source does not choose a field, topology/enrichment, continuous families of \(2\)-linear maps, equivalence of descent data, additive/group-completion procedure, spectrum or elliptic target. Thus no specific model or elliptic identification is asserted. The supplied summary says later Stolz--Teichner/Lurie frameworks realize related \(2\)-categorical motivation in spirit, while distinguishing them from a completed literal Kapranov--Voevodsky construction.
\subsection{Short English statement}
Investigate whether coherent \(2\)-vector-space transition data can construct elliptic cohomology, replacing strict bundle cocycles by natural isomorphisms with higher coherence.
\subsection{Sources}
\begin{itemize}
\item Mark Hovey, \emph{Algebraic Topology Problem List}, ``Elliptic cohomology,'' item 2, surviving mirror \url{https://www-users.cse.umn.edu/~tlawson/hovey/}. Recorded access 2026-07-29; extraction html-ordered-list.
\item Kapranov--Voevodsky, source-named paper on \(2\)-vector spaces, without a title or full bibliography entry. The summary names Stolz--Teichner and Lurie frameworks as related but not literal completion, without exact references.
\item Local UnsolvedMath registry, \texttt{data/unsolvedmath/problems.json}, permanent ID \(11100022\), exact source code \texttt{catalogue entry}; supplied partially\_solved/PARTIAL-PROGRESS, with partial/superseded qualification.
\end{itemize}
\end{document}
